\documentclass[10pt,a4paper]{article}
\usepackage[top=12mm,bottom=12mm,left=14mm,right=14mm,headsep=2mm,footskip=7mm]{geometry}
\usepackage{fontspec}
\usepackage{xeCJK}
\setmainfont{Noto Sans}
\setsansfont{Noto Sans}
\setCJKmainfont{Noto Sans CJK SC}
\setCJKsansfont{Noto Sans CJK SC}
\usepackage{xcolor}
\usepackage{graphicx}
\usepackage{tikz}
\usetikzlibrary{calc,arrows.meta,positioning}
\usepackage{fancyhdr}
\usepackage{hyperref}
\usepackage{ragged2e}
\usepackage{microtype}
\definecolor{P2Paper}{HTML}{FCFAF7}
\definecolor{P2Ink}{HTML}{2D3238}
\definecolor{P2Muted}{HTML}{70757A}
\definecolor{P2BlueLight}{HTML}{D3EAF5}
\definecolor{P2Apricot}{HTML}{F5D9B8}
\definecolor{P2Rose}{HTML}{E8C7D3}
\definecolor{P2Violet}{HTML}{D7D3EA}
\definecolor{P2Blue}{HTML}{5B9FBE}
\definecolor{P2ApricotAccent}{HTML}{D49757}
\definecolor{P2RoseAccent}{HTML}{B97589}
\definecolor{P2VioletAccent}{HTML}{8177A8}
\definecolor{P2Rule}{HTML}{DDE3E7}
\definecolor{P2Panel}{HTML}{FFFDFC}
\pagecolor{P2Paper}
\color{P2Ink}
\setlength{\parindent}{0pt}
\setlength{\parskip}{3.5pt}
\linespread{1.04}
\hypersetup{hidelinks}
\pagestyle{fancy}
\fancyhf{}
\renewcommand{\headrulewidth}{0pt}
\fancyfoot[L]{\fontsize{6.6}{8}\selectfont\color{P2Muted}PROCESS REPORT · WORKFLOW CONSTRUCTION}
\fancyfoot[R]{\fontsize{6.6}{8}\selectfont\color{P2Muted}\thepage}
\newcommand{\eyebrow}[1]{{\fontsize{7.8}{9.4}\selectfont\bfseries\color{P2RoseAccent}#1}\par\vspace{1pt}}
\newcommand{\pagetitle}[1]{{\fontsize{18.5}{21.5}\selectfont\bfseries\color{P2Ink}#1}\par\vspace{3pt}}
\newcommand{\deck}[1]{{\fontsize{9.0}{13.0}\selectfont\color{P2Ink}\RaggedRight #1}\par\vspace{4pt}}
\newcommand{\tracehead}[2]{{\fontsize{7.1}{8.4}\selectfont\bfseries\color{#1}#2}\par}
\newcommand{\tracetitle}[1]{{\fontsize{10.2}{12.2}\selectfont\bfseries\color{P2Ink}#1}\par\vspace{1.5pt}}
\newcommand{\tracebody}[1]{{\fontsize{8.0}{11.2}\selectfont\color{P2Ink}\RaggedRight #1}\par}
\newcommand{\provenance}[1]{\vspace{2pt}{\color{P2Rule}\hrule height .35pt}\vspace{2.5pt}{\fontsize{6.9}{9.2}\selectfont\color{P2Muted}\RaggedRight\textbf{Original Evidence.} #1}\par}
\tikzset{
 userA/.style={rounded corners=1mm,draw=P2ApricotAccent,line width=.55pt,fill=P2Apricot,fill opacity=.34},
 userR/.style={rounded corners=1mm,draw=P2RoseAccent,line width=.55pt,fill=P2Rose,fill opacity=.34},
 gptB/.style={rounded corners=1mm,draw=P2Blue,line width=.55pt,fill=P2BlueLight,fill opacity=.27},
 gptV/.style={rounded corners=1mm,draw=P2VioletAccent,line width=.55pt,fill=P2Violet,fill opacity=.27},
 arrA/.style={-{Stealth[length=2.0mm,width=1.45mm]},draw=P2RoseAccent,line width=.72pt},
 arrB/.style={-{Stealth[length=2.0mm,width=1.45mm]},draw=P2Blue,line width=.72pt},
 arrV/.style={-{Stealth[length=2.0mm,width=1.45mm]},draw=P2VioletAccent,line width=.72pt},
 tagA/.style={circle,minimum size=4.8mm,inner sep=0pt,draw=P2ApricotAccent,line width=.6pt,fill=P2Paper,font=\fontsize{6.7}{7}\selectfont\bfseries},
 tagR/.style={circle,minimum size=4.8mm,inner sep=0pt,draw=P2RoseAccent,line width=.6pt,fill=P2Paper,font=\fontsize{6.7}{7}\selectfont\bfseries},
 tagB/.style={circle,minimum size=4.8mm,inner sep=0pt,draw=P2Blue,line width=.6pt,fill=P2Paper,font=\fontsize{6.7}{7}\selectfont\bfseries}
}

\begin{document}

\eyebrow{WORKFLOW CONSTRUCTION · 26}
\pagetitle{流程图不是越多越好：先淘汰不适合研究叙事的视觉语言}
\deck{我们先试了多种Process Diagram。我没有因为“看起来高级”就全都保留：Metro/Storyline和Phase Ribbon直接淘汰，Swimlane我觉得太乱，Decision Tree则要求改成更适合A4横向阅读。}
\begin{tikzpicture}[x=1mm,y=1mm]
  \node[anchor=north west,inner sep=0] (u) at (0,0) {\includegraphics[width=43mm]{assets/user_reject_diagrams_phrase.png}};
  \node[anchor=north west,inner sep=0] (g) at (0,-31) {\includegraphics[width=52mm]{assets/gpt_diagram_filter.png}};
  \coordinate (ua) at ($(u.south east)+(-2mm,4mm)$);
  \coordinate (ga) at ($(g.north east)+(-4mm,-22mm)$);
  \draw[arrA] (ua) .. controls (80,-4) and (83,-42) .. (ga);
  \node[tagA] at ($(ua)+(3.5mm,0)$) {1};
  \node[anchor=north west,text width=109mm,align=left] at (69,-2) {
    \tracehead{P2ApricotAccent}{MICRO TRACE 32}\tracetitle{“1和5淘汰，2太乱，3横向，4要突出我的贡献”}
    \tracebody{我把五类图逐个看了一遍：哪些淘汰、哪些保留、哪些必须改。GPT随后也不再追求一个万能流程图，而是让不同图型各自回答不同问题。}
    \vspace{8mm}
    \tracehead{P2Blue}{SELECTION RULE}\tracetitle{图型由问题选择，不由复杂度选择}
    \tracebody{我后来判断流程图好不好，先看它到底回答什么：总体阶段、职责分工、候选淘汰、Evidence链还是收敛过程，而不是看它有多复杂。}
  };
\end{tikzpicture}
\provenance{我的原句与后续筛选结果来自同一轮Process Diagram Style Study复盘。}
\newpage

\eyebrow{WORKFLOW CONSTRUCTION · 27}
\pagetitle{Human Reasoning 必须成为流程图中心，而不是被 Evidence 抽象掉}
\deck{我对Evidence--Decision Graph最不满意的一点是：只画“证据 $\rightarrow$ 决策”，老师看不到我到底发现了什么、为什么觉得有问题、又提出了什么建议。这个反馈后来直接改变了最核心的Human-in-the-loop图。}
\begin{tikzpicture}[x=1mm,y=1mm]
  \node[anchor=north west,inner sep=0] (u) at (0,0) {\includegraphics[width=43mm]{assets/user_human_reasoning_phrase.png}};
  \node[anchor=north west,inner sep=0] (g) at (0,-33) {\includegraphics[width=50mm]{assets/gpt_human_reasoning.png}};
  \coordinate (ua) at ($(u.south east)+(-2mm,4mm)$);
  \coordinate (ga) at ($(g.north east)+(-4mm,-40mm)$);
  \draw[arrA] (ua) .. controls (80,-5) and (83,-60) .. (ga);
  \node[tagA] at ($(ua)+(3.5mm,0)$) {1};
  \node[anchor=north west,text width=111mm,align=left] at (67,-2) {
    \tracehead{P2ApricotAccent}{MICRO TRACE 33}\tracetitle{“要明确展示我发现、判断与建议”}
    \tracebody{我不是想在图上简单加一个“User”标签，而是要把Human reasoning真正展开：我发现了什么 $\rightarrow$ 为什么觉得有问题 $\rightarrow$ 我提出什么建议 $\rightarrow$ 后面怎么验证。}
    \vspace{8mm}
    \tracehead{P2VioletAccent}{CORE GRAMMAR}\tracetitle{GPT Proposal $\rightarrow$ Human Reasoning $\rightarrow$ Evidence / Experiment $\rightarrow$ Final Decision}
    \tracebody{这条链后来成为Process Report里最能看出我真正参与判断的一种图，也和Interaction Evidence的phrase-level标注保持一致。}
  };
\end{tikzpicture}
\provenance{该页使用我对方案4的真实批评及GPT随后形成的Human Reasoning--Evidence--Decision结构。}
\newpage

\eyebrow{WORKFLOW CONSTRUCTION · 28}
\pagetitle{最终收敛为三类 Process Diagram Language，而不是一个万能流程图}
\deck{几轮比较以后，我明确选了2和4作为主力，3保留给多候选场景，1降级，5淘汰。GPT随后把这次选择整理成三类稳定的Diagram Grammar。}
\begin{tikzpicture}[x=1mm,y=1mm]
  \node[anchor=north west,inner sep=0] (u) at (0,0) {\includegraphics[width=38mm]{assets/user_diagram_selection.png}};
  \node[anchor=north west,inner sep=0] (g) at (0,-23) {\includegraphics[width=48mm]{assets/gpt_final_diagram_language.png}};
  \coordinate (ua) at ($(u.south east)+(-2mm,3mm)$);
  \coordinate (ga) at ($(g.north east)+(-4mm,-26mm)$);
  \draw[arrA] (ua) .. controls (75,-3) and (80,-33) .. (ga);
  \node[tagA] at ($(ua)+(3.5mm,0)$) {1};
  \node[anchor=north west,text width=112mm,align=left] at (65,-2) {
    \tracehead{P2ApricotAccent}{MICRO TRACE 34}\tracetitle{最终选择由我明确给出}
    \tracebody{最后的选择是我自己给出的：2和4作为核心，3用于多候选场景，1降级，5淘汰。GPT负责把这个选择整理成后续可以稳定复用的图示语言。}
    \vspace{7mm}
    \tracehead{P2Blue}{TYPE A}\tracetitle{Editorial Decision Board}
    \tracebody{一般过程、阶段性结论、单页“背景—决策—结果”。}
    \vspace{5mm}
    \tracehead{P2VioletAccent}{TYPE B}\tracetitle{Human Reasoning $\rightarrow$ Evidence $\rightarrow$ Decision}
    \tracebody{核心Human-in-the-loop、AI批判、参数纠正与实验验证。}
    \vspace{5mm}
    \tracehead{P2RoseAccent}{TYPE C}\tracetitle{Horizontal Candidate Decision Tree}
    \tracebody{多候选、筛选、淘汰、对照实验和最终选择。}
  };
\end{tikzpicture}
\provenance{上方为我的最终选择；下方为GPT随后冻结的Process Diagram Language vFinal。}
\end{document}
